\section{Discussions}
\paragraph{Broad applicability of {\gymname}.}  
{\gymname} is implemented as a standard Gym environment, offering several advantages: (1) it exposes a familiar API (e.g., \textit{reset}, \textit{step}), allowing seamless integration with other Gym-compatible systems; and (2) it abstracts user behavior as a black-box environment, enabling agents to interact directly without requiring internal knowledge of user logic. This setup mirrors real-world scenarios where agents must infer user intent without full transparency, reducing the model's cognitive burden and encouraging generalizable behavior.

These traits also make {\gymname} friendly for both evaluation and training:  
(1) As a \textbf{benchmark}, it enables flexible testing configurations, including customizable option sets, reward functions (e.g., answer correctness, search and action validity, penalties, etc.), and user feedback mechanisms (e.g., rule-based, user simulated, etc.). This allows for fine-grained agent behavioral analysis under varied user profiles.  
(2) As a \textbf{training environment}, {\gymname} can be used by any model with tool use capabilities via a standardized interaction interface we provide. It supports both supervised fine-tuning and reinforcement learning, especially multi-turn RL, by providing turn-wise rewards and partial credit signals, which is critical for improving robustness in user-aligned interaction. This also makes {\gymname} particularly well-suited to the emerging trend of applying RL to agentic LLMs~\citep{jin2025search, li2025torl, qian2025toolrl}.

\paragraph{Balancing efficiency and effectiveness in user interaction.}
Our weighted timing analysis reveals that many models, including GPT-4o, often prioritize efficiency by guessing answers early rather than thoroughly probing user preferences. As a result, while agents may appear helpful by providing quick responses, they frequently overlook deeper user intent. This mirrors real-world experiences where users receive seemingly helpful answers but still feel misunderstood, leading to repeated prompt revisions and dissatisfaction.

Conversely, we rarely observed the opposite failure mode, where models hesitating indefinitely and over-requesting information. Yet both extremes underscore the need to balance efficiency (responding promptly) and effectiveness (satisfying user needs). This trade-off also reflects a broader concern around tool use efficiency, which has received increasing attention in recent works~\citep{qian2025smart, wang2025acting}. While {\gymname} does not explicitly quantify this balance, it surfaces model behaviors that targeted training can effectively address.
Future work, especially using RL, can shape agent behavior toward this balance. To promote efficiency, reward functions can penalize late preference discovery by decaying rewards based on the turn when the correct answer is given. To promote effectiveness, partial rewards can be granted for each user preference successfully elicited, and penalties applied to correct answers that lack sufficient supporting interaction. These strategies discourage guessing and encourage deliberate, user-aware interactions. {\gymname} supports such reward customizations, providing a flexible platform for training agents that are not only capable, but also user-aligned.
